\section{\gls{CFD}\label{ch3:cfdBackround}}
        The roots of \gls{CFD} can be traced back to the $1950$s and $1960$s when
        researchers began using computers to solve \glspl{PDE} that describe
        fluid flow. This period saw the initial application of numerical
        analysis to solve simplified fluid problems.\\

        Numerical analysis provides methods and algorithms that are suitable
        for \gls{CFD}, such as \gls{FDM}, \gls{FVM}, and \gls{FEM}
        \cite{Greenshields2022}.
        \begin{itemize}
            \item \textbf{\gls{FDM}:} Approximates derivatives, usually using
                Taylor series expansions. These methods are simpler to
                implement and well-suited for structured grids. However, these
                methods perform poorly for complex flows with strong gradients,
                nonequilibrium properties, or complex geometries.
            \item \textbf{\gls{FVM}:} Integrates the governing equations over a
                control volume and applies the divergence theorem to convert
                volume integrals into surface integrals. It enforces local
                conservation of fluxes across volume boundaries. This method is
                applicable to unstructured meshes and well-suited for complex
                flows with strong conservation requirements.
            \item \textbf{\gls{FEM}:} Formulates the \gls{PDE} as a variational
                problem (typically a \textit{weak} or \textit{Galerkin} form).
                The solution is approximated using basis functions defined over
                elements. This method is very flexible in handling complex
                geometries and boundary conditions, although it is more
                mathematically involved.  
        \end{itemize}

        In the $1970$s and $1980$s, the field of \gls{CFD} has been established as a
        distinct discipline, and the availability of more powerful computers
        allowed researches to address more complex and realistic problems. This
        development was driven by the aerospace community \cite{Anderson1995}.\\

        Between the $1980$s to $1990$s, \gls{CFD} became a more widely and accessible
        with development of commercial \gls{CFD} software such as ANSYS Fluent.
        During this period, the ability to simulate and analyze fluid flows
        using computational methods gained popularity in industries such as
        aerospace and automotive engineering. From the $1990$s to the present,
        continuous advancements in numerical
        algorithms, grid generation techniques, and parallel computing have been
        made. Adaptive mesh refinement, unstructured grids, and
        \gls{HPC} have significantly improved the
        accuracy and efficiency of \gls{CFD} simulations.\\

        In recent years, there has been a push towards multidisciplinary simulations,
        integrating \gls{CFD} with other disciplines such as structural
        mechanics, heat transfer and optics. This integration allows for solving
        more realistic and complex problems.
